\bookchapter{Launch Files and Parameters}\label{ch:ros3}

\section*{Explanation}
A real robot runs ten to fifty nodes. Starting them by hand in ten terminals, each with the right topic remappings and settings, is slow and error-prone: one forgotten \verb!--ros-args! and the robot does not see the ball. A \textbf{launch file} describes the whole system once: which nodes, with which names, namespaces, remappings and parameters. \textbf{Parameters} are the settings of a node (thresholds, topic names, the team colour), kept outside the code so that you can change them on the field in a minute without rebuilding. This chapter shows both, and how to organise a package so that one command starts your robot.

\section*{Parameters}
A parameter has a name, a type (bool, integer, double, string, or an array of these) and a value. A node must \textbf{declare} the parameters it uses, with a default value; undeclared parameters cannot be set (a good thing: a typo in the config file is caught).
\begin{verbatim}
import rclpy
from rclpy.node import Node

class BallDetector(Node):
    def __init__(self):
        super().__init__("ball_detector")
        self.declare_parameter("hue_min", 5)
        self.declare_parameter("hue_max", 20)
        self.declare_parameter("min_radius_px", 6.0)
        self.declare_parameter("team_colour", "blue")
        self.hue_min = self.get_parameter("hue_min").value
        self.get_logger().info(f"hue range starts at {self.hue_min}")
\end{verbatim}
The command line shows and changes them while the node runs:
\begin{verbatim}
$ ros2 param list /ball_detector
$ ros2 param get /ball_detector hue_min
$ ros2 param set /ball_detector hue_min 8
$ ros2 param dump /ball_detector > detector.yaml   # save current values
\end{verbatim}
A \verb!set! changes only the value in the running node. If the code reads the value once in \verb!__init__! (as above), nothing happens. To react to changes register a callback, which can also \emph{reject} wrong values (for example \verb!hue_min! larger than \verb!hue_max!):
\begin{verbatim}
from rcl_interfaces.msg import SetParametersResult
...
        self.add_on_set_parameters_callback(self.on_params)

    def on_params(self, params):
        for p in params:
            if p.name == "hue_min":
                self.hue_min = p.value
        return SetParametersResult(successful=True)
\end{verbatim}
\begin{keyidea}[Parameters are for settings and not for data]
Parameters change rarely (thresholds, filenames, gains, topic names). Anything that changes continuously, such as a target position, is a topic. Anything that is an \emph{event} with an answer (``recalibrate now'') is a service. Mixing these up gives programs that are hard to understand and to test.
\end{keyidea}

\subsection*{Parameter files in YAML}
The usual way to give parameters to nodes is a YAML file. The first key is the node name (or \verb!/**! for every node), the second is always \verb!ros__parameters! (two underscores):
\begin{verbatim}
# config/detector.yaml
ball_detector:
  ros__parameters:
    hue_min: 5
    hue_max: 20
    min_radius_px: 6.0
    team_colour: "blue"
"/**":                        # applies to every node of this launch
  ros__parameters:
    use_sim_time: false
\end{verbatim}
A common beginner error: an integer where the code declares a double (\verb!6! instead of \verb!6.0!) leads to a type error, because the declared default fixes the type. Use one YAML file for what is the same on all robots, and a second for what differs per robot (camera calibration, joint offsets); the later file wins.

\section*{Launch files}
ROS~2 launch files are normally written in \textbf{Python} (the examples here), but XML and YAML formats exist and are fine for simple cases; the Python form is the most powerful because you have the whole language, for example to choose nodes by the hostname. A launch file defines a function \verb!generate_launch_description! that returns a list of actions.
\begin{verbatim}
# launch/vision.launch.py
import os
from ament_index_python.packages import get_package_share_directory
from launch import LaunchDescription
from launch.actions import DeclareLaunchArgument
from launch.substitutions import LaunchConfiguration
from launch_ros.actions import Node

def generate_launch_description():
    share = get_package_share_directory("my_robot_vision")
    cfg = os.path.join(share, "config", "detector.yaml")
    return LaunchDescription([
        DeclareLaunchArgument("camera_topic",
                              default_value="/camera/image_raw"),
        DeclareLaunchArgument("use_sim_time", default_value="false"),
        Node(
            package="my_robot_vision",
            executable="ball_detector",
            name="ball_detector",
            namespace="robot1",
            parameters=[cfg, {"use_sim_time":
                              LaunchConfiguration("use_sim_time")}],
            remappings=[("image",
                         LaunchConfiguration("camera_topic"))],
            output="screen",
        ),
    ])
\end{verbatim}
Start it with \verb!ros2 launch my_robot_vision vision.launch.py use_sim_time:=true!. Reading the example:
\begin{itemize}[nosep]
\item The action \verb!DeclareLaunchArgument! defines a command-line argument with a default. \verb!LaunchConfiguration("name")! is a \emph{substitution}: a placeholder evaluated when the launch starts, not when the file is read (do not use it as an ordinary Python string).
\item \texttt{Node(\ldots)} starts one executable. \verb!parameters! takes files and dictionaries; later entries win. \verb!remappings! is a list of (old, new) pairs.
\item \verb!output="screen"! prints the node's output in the terminal.
\item \verb!ros2 launch ... --show-args! lists the arguments of a launch file.
\end{itemize}
\subsection*{Building bigger systems}
\begin{verbatim}
from launch.actions import IncludeLaunchDescription, GroupAction
from launch.launch_description_sources import PythonLaunchDescriptionSource
from launch.conditions import IfCondition
from launch_ros.actions import PushRosNamespace

camera = IncludeLaunchDescription(
    PythonLaunchDescriptionSource(
        os.path.join(cam_share, "launch", "camera.launch.py")),
    launch_arguments={"fps": "30"}.items())

robot = GroupAction([PushRosNamespace("robot1"),
                     camera, detector, behaviour])

viewer = Node(package="rviz2", executable="rviz2",
              condition=IfCondition(LaunchConfiguration("gui")))
\end{verbatim}
\verb!IncludeLaunchDescription! reuses other launch files (the camera package ships its own), \verb!GroupAction! with \verb!PushRosNamespace! puts everything inside a namespace (a second robot: the same launch file with \verb!robot2!), and a \verb!condition! switches nodes on and off with an argument, for example no \verb!rviz2! on the robot. In Python you can also write ordinary logic around it:
\begin{verbatim}
import socket
is_robot = socket.gethostname().startswith("k1-")
actions = [camera, detector] + ([] if is_robot else [viewer])
\end{verbatim}
Other useful actions: \verb!ExecuteProcess! (start any command), \verb!TimerAction! (start something after $5$\,s, for example the behaviour when the camera is up), \verb!RegisterEventHandler! with \verb!OnProcessExit! (react when a process ends: restart, or stop everything), and the \verb!respawn=True! option of a \verb!Node! for a crashing but unimportant node. For a safety-relevant node, respawning is \emph{not} the answer; see Chapter~\ref{ch:hw8}.

\section*{Package layout}
Launch and parameter files are \emph{installed} with the package. For a Python package the folder looks like this:
\begin{verbatim}
my_robot_vision/
  package.xml
  setup.py
  my_robot_vision/ball_detector.py
  launch/vision.launch.py
  config/detector.yaml
\end{verbatim}
and \verb!setup.py! must copy the two folders into the \verb!install! directory, otherwise the package folder lookup in the launch file finds nothing:
\begin{verbatim}
import os
from glob import glob
from setuptools import setup
...
    data_files=[
        ("share/ament_index/resource_index/packages",
         ["resource/my_robot_vision"]),
        ("share/my_robot_vision", ["package.xml"]),
        (os.path.join("share", "my_robot_vision", "launch"),
         glob("launch/*.launch.py")),
        (os.path.join("share", "my_robot_vision", "config"),
         glob("config/*.yaml")),
    ],
\end{verbatim}
With a C++ package, \verb!CMakeLists.txt! contains an \verb!install(DIRECTORY launch config ...)! command for the two folders. After editing, rebuild (\verb!colcon build --symlink-install! lets changes in the launch and YAML files take effect without another build) and source again.

\section*{Starting at boot}
On the robot you want the system to start without a laptop. A \verb!systemd! service (Chapter~\ref{ch:linux5}) runs a small wrapper script that sources the environment, then launches:
\begin{verbatim}
#!/bin/bash
# /opt/robot/start.sh
source /opt/ros/lyrical/setup.bash
source /home/robot/ros2_ws/install/setup.bash
export ROS_DOMAIN_ID=17
exec ros2 launch my_robot_bringup robot.launch.py
\end{verbatim}
The service file contains \verb!ExecStart=/opt/robot/start.sh!, \verb!Restart=on-failure! and \verb!User=robot!; \verb!journalctl -u robot -f! shows the output. The \verb!exec! makes the launch the main process, so that \verb!systemctl stop! reaches it and it can shut the nodes down cleanly.

\subsection*{Composition in one sentence}
Several nodes can be loaded into \emph{one process} (a \emph{component container}): they then pass messages by pointer without copying and without the network, which matters for big images. In a launch file this is the \verb!ComposableNodeContainer!. Use it when a profile shows that copying images costs too much (Chapter~\ref{ch:py12} shows how to measure), not before.

\section*{Practice problems}
\begin{enumerate}[label=\textbf{\thechapter.\arabic*}]
  \item Write a YAML parameter file for the node \verb!walk_controller! with the parameters \verb!step_height! ($0.04$), \verb!step_period! ($0.5$), \verb!max_speed! ($0.4$) and \verb!gait! (\verb!"trot"!).
  \item The file of problem 1 says \verb!max_speed: 1! while the node declares \verb!declare_parameter("max_speed", 0.4)!. What happens, and why?
  \item Write the command that starts a launch file \verb!robot.launch.py! of the package \verb!my_robot_bringup! with the arguments \verb!team:=red! and \verb!gui:=false!.
  \item In a launch file you write \verb!if LaunchConfiguration("gui") == "true":!. Why does that not work, and what do you use instead?
  \item A node reads its \verb!hue_min! parameter once in \verb!__init__!. You run \verb!ros2 param set!; the value changes in \verb!ros2 param get!, but not the behaviour. Explain and give two ways to fix it.
  \item Explain what \verb!namespace="robot1"! and \verb!remappings=[("image", "/cam/front")]! do to the topic \verb!image! that the node subscribes to. What is its full name in each case, and when is the remapping applied?
  \item Your launch file works from \verb!ros2 launch! in the source folder, but after installing it says \verb!package 'my_robot_vision' not found! or cannot find \verb!detector.yaml!. List three things to check.
  \item You want the behaviour node to start only $3$\,s after the camera. Which launch action do you use? Why is a fixed delay not ideal, and what would be more robust?
  \item Why should you not use \verb!respawn=True! for the node that sends the walking command? What should happen instead if it crashes?
  \item[\textbf{\thechapter.10}$\star$] Write a complete launch file for a two-robot simulation: for each of \verb!robot1! and \verb!robot2! it includes \verb!vision.launch.py! inside the namespace, with the parameter \verb!team_colour! red for the first and blue for the second, and one shared \verb!rviz2! only if the argument \verb!gui! is true.
\end{enumerate}
